% src: RESULTS.md header (341 runs, 1756 cases, identical protocol);
% src: RESULTS_v1_vs_v2.md footer (main 195, curves 72, corrupt 36, backend 20, twod 18) — verified against runs_v2.csv stage counts;
% src: RESULTS.md §1.0б (R² with a single global denominator; R²_n over normal components);
% src: publication/code/trainer.py line 442 (macro_rmse = mean of the four component RMSEs); publication/code/stats.py line 147 (scipy ttest_rel over seeds).
All seed-replicated results below were obtained on the cleaned dataset of 1756 cases (Section~\ref{sec:prep-data}) in 341 training runs: 195 for the split matrix, 72 for the learning curves, 36 for label corruption, 20 for the backend comparison and 18 for the two-dimensional ablation; the one exception, the single-run comparison of Table~\ref{tab:admissibility}, is described in Section~\ref{sec:res_fem}. % EDIT: "all results ... 1756" qualified — tab:admissibility uses the 1498/261 run (critic; ERRATA.md E-24)
The three one-dimensional families -- the data-driven MLP, the strong-form PINN and the weak-form VPINN -- share architecture, optimiser, budget, early stopping and seeds and differ only in the composition of the loss (Section~\ref{sec:models}). Accuracy is reported as macro-RMSE, the mean of the four component RMSEs in MPa, as mean~$\pm$~standard deviation over five seeds; differences between families are tested with a paired $t$-test over seeds \citep{demsar2006} and called established at $p<0.05$. Because $\tau_{rz}$ is two orders of magnitude smaller than the normal components, a coefficient of determination over the normal components only, $R^2_{\mathrm{n}}$, with a single denominator (global denominator: the variance of each component over the whole dataset) is given alongside; where a per-component $R^2$ or NRMSE below uses the held-out slice as denominator, this is stated at the number. % EDIT: denominator named (critic)
The finite-element model is the reference throughout; no physical measurement enters, so the section characterises the surrogates relative to their data source, not to the process (see the closing paragraph of this section). % EDIT: \ref{sec:limitations} replaced by plain text — no such section is written; a Limitations paragraph closes Section~\ref{sec:res_fem} (critic)

% =====================================================================
\subsubsection{Extrapolation versus interpolation}
\label{sec:res_extrap}

% src: RESULTS.md §1 (macro-RMSE table, n_test, 5 seeds) — all 39 entries verified; RESULTS.md §1.0б (constant predictor, all 13 values verified);
% src: publication/code/splits.py REGIONS (alpha=20, alpha=12, Q=0.20, Q=0.10, k=1.0, v=250, corner = alpha=20 & Q>=0.15) and make_random_split frac=0.15 (0.15 × 1756 = 263);
% src: PLAYBOOK.md §7.1 (matched = random hold-out of the same size; corner is a hole in the joint space, not out-of-range on any single axis);
% src: ERRATA.md E-22 (Q is the reduction in diameter).
% EDIT: the macro-RMSE table (a third \label{tab:splits}, duplicate of T1) deleted — see Table~\ref{tab:main} in tables/all_tables.tex. Its "Constant" column survives in Table~\ref{tab:axis} (five regions) and in the text (random: 66.0 MPa).

% src: RESULTS.md §1 and §1.0б (random 8.0 MPa, R²_n = +0.99 all families; constant predictor 66.0 MPa).
Table~\ref{tab:main} gives the macro-RMSE on all thirteen splits. On the random hold-out all three families reach 8.0~MPa with $R^2_{\mathrm{n}}=0.99$ (global denominator), against 66~MPa for the constant predictor, and are indistinguishable. The informative contrast is between each region split and its matched control of the same size, which separates the cost of leaving the training range from the cost of having fewer training cases.

% src: RESULTS.md §1.0б "Цена экстраполяции зависит от оси" table (family-averaged RMSE outside / control / ratio / R²_n / null model) — all 25 entries verified.
\begin{table*}[t]
\caption{Cost of holding out a whole factor level, averaged over the three families: macro-RMSE on the held-out region, on the matched control of the same size, their ratio, $R^2_{\mathrm{n}}$ (global denominator) on the held-out region, and the constant predictor (MPa).} % EDIT: denominator named (critic)
\label{tab:axis}
\centering\small
\begin{tabular}{@{}lrrrrr@{}}
\toprule
Held-out region & Outside & Control & Ratio & $R^2_{\mathrm{n}}$ outside & Constant \\
\midrule
$k=1.0$                              & 8.9  & 8.3 & $\times1.08$ & $+0.98$ & 67 \\
$\alpha=20^\circ$                    & 13.2 & 8.1 & $\times1.64$ & $+0.96$ & 61 \\
$\alpha=20^\circ$, $Q\ge0.15$        & 16.1 & 7.3 & $\times2.20$ & $+0.94$ & 61 \\
$Q=0.20$                             & 24.7 & 7.8 & $\times3.15$ & $+0.81$ & 60 \\
$v=250$~m/min                        & 50.9 & 8.6 & $\times5.93$ & $+0.36$ & 71 \\
\bottomrule
\end{tabular}
\end{table*}

% src: RESULTS.md §1.0б (ratios ×1.08, ×1.64, ×2.20, ×3.15, ×5.93; R²_n on Q_high: MLP +0.78, PINN +0.81, VPINN +0.83);
% src: RESULTS.md §1.1 (ΔRMSE vs control: MLP k_max 1.13 ± 0.58 established; PINN k_max 0.11 ± 0.40 not established).
The cost of extrapolation depends on the axis by almost an order of magnitude (Table~\ref{tab:axis}). Holding out the longest die land ($k=1.0$) costs practically nothing: the excess over the control is $1.13\pm0.58$~MPa for the MLP (established by Welch's test, Section~\ref{sec:models}) and $0.11\pm0.40$~MPa for the PINN (not established). % EDIT: "within the seed spread" held only for the PINN; RESULTS.md §1.1 marks MLP/k_max as established (critic)
Holding out the largest semi-angle costs a factor of 1.6 with $R^2_{\mathrm{n}}$ (global denominator) still at 0.96, the largest reduction a factor of 3.2 ($R^2_{\mathrm{n}}=0.78$--$0.83$), the highest speed a factor of 5.9. The spread between axes is far larger than any difference between families, so a single ``extrapolation error'' would be uninformative.

% src: RESULTS.md §1 (interp:alpha_mid 8.69–9.06 vs extrap:alpha_max 12.98–13.63 vs random 7.98–8.04; interp:Q_mid 10.96–12.40 vs extrap:Q_high 23.41–26.27);
% src: RESULTS.md §1.0б per family: alpha_mid R²_n +0.98/+0.98/+0.98; alpha_max +0.96 (MLP) / +0.97 (PINN) / +0.97 (VPINN) — the family average +0.96 is what tab:axis quotes.
% NOTE: RESULTS.md text quotes "10.4 MPa" for the interior alpha level; that figure is the mean over BOTH interior splits (alpha_mid 8.91 and Q_mid 11.85 → 10.38) — a text-generation slip in make_results.py; the table values are used here.
Removing an interior level is not free either. Along $\alpha$, where the volumes are nearly equal, withholding the interior level $\alpha=12^\circ$ gives 8.7--9.1~MPa ($R^2_{\mathrm{n}}=0.98$, global denominator) and withholding the extreme level $\alpha=20^\circ$ gives 13.0--13.6~MPa ($R^2_{\mathrm{n}}=0.96$--$0.97$), against 8.0~MPa on the random hold-out: most of the price is paid for the hole in a five-level factor grid, and the direction of the hole adds comparatively little. Along $Q$ the same pattern is steeper, 11.0--12.4~MPa for $Q=0.10$ against 23.4--26.3~MPa for $Q=0.20$.

% src: RESULTS.md §1.0б (v_max: 48–56 vs 71.4 constant; R²_n +0.23 / +0.45 / +0.40); RESULTS.md §1.3б (τ_rz NRMSE 1.622 / 1.945 / 1.879 on v_max); RESULTS.md §1 (MLP sd 19.04); VALIDATION.md §6 (speed level coincides with simulation batch).
The highest speed is where the surrogate breaks down. With $v=250$~m/min withheld the error is 48--56~MPa, still below the 71~MPa of the constant predictor, but $R^2_{\mathrm{n}}$ (global denominator) falls to 0.23--0.45, the shear component is predicted worse than a constant (NRMSE 1.6--1.9, held-out standard deviation as denominator) and % EDIT: denominators named (critic)
the MLP becomes unstable across seeds ($\pm19$~MPa). This axis carries a further caveat: in the source simulations each speed level coincides with a separate simulation batch (Section~\ref{sec:prep-data}), so the speed factor partly encodes batch identity, and no claim about the physical effect of speed is made from this split.

% src: RESULTS.md §1.3 (paired t-tests over seeds, all 13 splits; sign flipped to PINN − MLP and VPINN − MLP) — all 26 Δ/p pairs verified; RESULTS_v1_vs_v2.md §2 (four significant splits).
% EDIT: the paired-test table (tab:tests, duplicate of T2) deleted — see Table~\ref{tab:paired} in tables/all_tables.tex, which also carries the t-statistics and the Holm correction. Its Δ(VPINN−MLP) on matched:Q_high read −0.23; T2 (and the csv, −0.2248) give −0.22.

% src: RESULTS.md §1.3 (all 13 mlp−pinn differences positive; p-values as in table; pinn−vpinn established on extrap:k_max Δ = −0.695, p = 0.018, AND on matched:v_max Δ = −0.111, p = 0.046 with |Δ| below the seed σ = 0.126);
% src: RESULTS.md §1.2 (Friedman p = 0.04076 extrap:alpha_max, 0.02237 extrap:k_max, 0.006738 matched:v_max; extrap_joint:corner p = 0.09072; Nemenyi separates only MLP–PINN; non-transitivity, ERRATA E-15);
% src: RESULTS_v1_vs_v2.md §3 and §8 (differences 0.3–1.9 MPa at a level of 8–16 MPa).
Table~\ref{tab:paired} summarises the seed-paired tests. The PINN has a lower macro-RMSE than the MLP on all thirteen splits, and the difference is established on four: the semi-angle extrapolation ($\Delta=-0.58$~MPa, $p=0.040$), the land-length extrapolation ($-1.16$~MPa, $p=0.033$), the joint corner ($-1.06$~MPa, $p=0.011$) and the matched control of the speed split ($-0.29$~MPa, $p=0.006$). With a Holm correction over the thirteen splits none of the pairwise differences remains significant (smallest adjusted $p = 0.079$ for PINN$-$MLP, 0.195 for VPINN$-$MLP); the four verdicts are read as a consistent direction, not as isolated proofs. % src: tables/all_tables.tex tab:paired caption (paper_tables.py, Holm over 13 splits: smallest adjusted p = 0.079 / 0.195) % EDIT: Holm sentence added (critic)
The VPINN is established better than the MLP on two splits (\texttt{interp:Q\_mid}, $p=0.015$; \texttt{matched:v\_max}, $p=0.022$). The strong form is established better than the weak form on \texttt{extrap:k\_max} ($-0.70$~MPa, $p=0.018$) and, by a margin smaller than the seed spread, on \texttt{matched:v\_max} ($-0.11$~MPa, $p=0.046$); on the remaining eleven splits the two physics-informed forms are indistinguishable. A Friedman test over seeds \citep{friedman1937,demsar2006} rejects equality on three of the four splits where the paired test does (\texttt{extrap:alpha\_max}, $p=0.041$; \texttt{extrap:k\_max}, $p=0.022$; \texttt{matched:v\_max}, $p=0.007$; on the corner $p=0.091$), with the Nemenyi post-hoc separating only the MLP--PINN pair; since indistinguishability there is not transitive, no ``single cluster'' is claimed. In absolute terms the advantage is small, 0.3--1.9~MPa at an error level of 8--16~MPa and of the same order as the seed spread; accuracy alone does not justify the physics term (Section~\ref{sec:res_fem}).

% src: RESULTS.md §3 (learning curves, macro-RMSE at 10/25/50/100 %) — all 24 entries verified; RESULTS_v1_vs_v2.md §5 (gaps, slopes on log fraction, p-values);
% src: runs_v2.csv stage=curves (n_train 149/373/746/1493 random, 138/344/688/1377 alpha_max; 3 seeds); paired p-values recomputed from the csv: 0.317/0.099/0.048/0.718 and 0.498/0.033/0.886/0.238; slopes +0.197 (p = 0.211) and −0.355 (p = 0.070).
\begin{table*}[t]
\caption{Learning curves: macro-RMSE (MPa) versus the fraction of the training pool, three seeds. $\Delta$ is PINN minus MLP with the seed-paired $p$-value. Slope of $\Delta$ on $\log(\text{fraction})$: $+0.20$~MPa ($p=0.21$) on \texttt{random}, $-0.35$~MPa ($p=0.070$) on \texttt{extrap:alpha\_max}.}
\label{tab:curves}
\centering\small
\begin{tabular}{@{}lrrrrrrr@{}}
\toprule
Split & Fraction & $n_{\mathrm{train}}$ & MLP & PINN & VPINN & $\Delta$ & $p$ \\
\midrule
\texttt{random} & 10\,\%  & 149  & $13.09\pm0.67$ & $12.48\pm0.18$ & $12.98\pm0.10$ & $-0.61$ & 0.317 \\
                & 25\,\%  & 373  & $9.92\pm0.40$  & $9.46\pm0.39$  & $9.92\pm0.18$  & $-0.46$ & 0.099 \\
                & 50\,\%  & 746  & $8.80\pm0.15$  & $8.34\pm0.06$  & $8.53\pm0.21$  & $-0.46$ & 0.048 \\
                & 100\,\% & 1493 & $8.19\pm0.34$  & $8.09\pm0.20$  & $8.03\pm0.19$  & $-0.10$ & 0.718 \\
\addlinespace
\texttt{extrap:alpha\_max} & 10\,\%  & 138  & $20.48\pm0.63$ & $20.72\pm0.21$ & $20.68\pm0.67$ & $+0.24$ & 0.498 \\
                           & 25\,\%  & 344  & $16.74\pm0.99$ & $17.33\pm1.17$ & $17.03\pm0.33$ & $+0.59$ & 0.033 \\
                           & 50\,\%  & 688  & $14.01\pm0.40$ & $14.06\pm0.70$ & $14.17\pm0.95$ & $+0.05$ & 0.886 \\
                           & 100\,\% & 1377 & $13.48\pm0.58$ & $12.95\pm0.36$ & $12.70\pm0.48$ & $-0.54$ & 0.238 \\
\bottomrule
\end{tabular}
\end{table*}

% src: RESULTS_v1_vs_v2.md §5 (opposite signs; slopes +0.197 p = 0.211, −0.355 p = 0.070; isolated significances at random 50 % and alpha_max 25 %; thesis not supported);
% src: RESULTS.md §3 quotes p = 0.255 at 10 % random: that is an unpaired Welch test (recomputed 0.2549 from the csv); the seed-paired test used throughout gives 0.317.
The learning curves (Table~\ref{tab:curves}) do not support the expectation that a physics prior pays off mainly when data are scarce. On the random split the PINN--MLP gap shrinks with more data ($-0.61\to-0.10$~MPa), as that expectation predicts; on the semi-angle extrapolation it moves the other way ($+0.24\to-0.54$~MPa). Neither trend is established (slopes $+0.20$~MPa per log-unit, $p=0.21$, and $-0.35$~MPa, $p=0.070$); the two isolated significant points ($p=0.048$ at 50\,\%, $p=0.033$ at 25\,\%) have opposite signs, as eight comparisons produce by chance, and at the smallest fraction the gap is not established on either split. The data-scarcity claim is therefore not made.

% =====================================================================
\subsubsection{Backend comparison}
\label{sec:res_backend}

% src: RESULTS.md §2 and ERRATA.md E-19 (derivative agreement 8.7e-15 abs, 9.8e-16 rel, jax_twin.py::check_gradient_agreement; DeepXDE configured with DDE_BACKEND=pytorch → two engines, not three);
% src: RESULTS.md §2 (implementations share everything but draw initial weights and batch order from their own RNGs: torch.manual_seed vs numpy default_rng).
The strong-form PINN was trained in two independent implementations of the same protocol, in PyTorch and in JAX, on the random split and the semi-angle extrapolation, five seeds each. The engines were checked first: on bit-identical weights and inputs the derivative $\partial\sigma_{rr}/\partial r$ from \texttt{torch.autograd} and from \texttt{jax.grad} agrees to $8.7\times10^{-15}$ absolute ($9.8\times10^{-16}$ relative), i.e.\ to machine precision. A DeepXDE implementation \citep{lu2021} of the same network runs on the PyTorch backend and differentiates with the same \texttt{torch.autograd.grad}, so the comparison involves two automatic-differentiation engines, not three.
% REV K22: reference to an earlier version of the text removed. The two implementations share architecture, data, split, nominal hyperparameters and stopping rule; the JAX twin applies $\lambda_{\mathrm{phys}} = 0.3$ directly, without the gradient-norm calibration of Eq.~\eqref{eq:calib} (effective weight 0.3 against $\approx 0.05$ in PyTorch), and draws initial weights and batch order from its own random stream. % src: jax_twin.py JaxTrainer.fit (no _calibrate; lam_p = 0.3); runs_v2.csv stage backend (phys_gain = 1.0 in all 10 jax rows, 0.132–0.215 in the 10 torch rows; 0.3 × 0.165 ≈ 0.05) % EDIT: the calibration confound named explicitly (critic; §2.3 "Backends"). The caption of tab:backend (paper_tables.py) carries the same wording.

% src: RESULTS_v1_vs_v2.md §7 (JAX 8.12 ± 0.18 vs PyTorch 7.98 ± 0.21, Δ +0.14, p = 0.269; 13.49 ± 0.30 vs 13.05 ± 0.49, Δ +0.43, p = 0.169;
%      full training 313 ± 102 vs 516 ± 94 s, ×1.65; per epoch 0.575 vs 0.745 s, ×1.30, p = 0.0009; epochs 538 vs 695, ×1.29); all recomputed from runs_v2.csv stage=backend (per-epoch p = 0.00086, unpaired).
% EDIT: the p = 0.0009 of RESULTS_v1_vs_v2.md §7 is an UNPAIRED equal-variance t-test (Welch 0.0010); the protocol declares the seed as the block, so the text uses the paired value of T4, p = 0.005 (n = 10, paired on split × seed); wall time p = 0.002, epochs p = 0.021 (tables/all_tables.tex tab:backend).
% src: runs_v2.csv (n_test 263 and 379; max epochs = 800, reached by jax/random seed 1, jax/alpha_max seed 0, torch/alpha_max seed 2); trainer.py line 172 / jax_twin.py line 40 (float64, jax_enable_x64); run_matrix.py lines 87/129 (torch.set_num_threads(1); JAX thread count not pinned — TODO state in §2.3).
% NOTE: RESULTS.md §2 marks the alpha_max gap as "established" under its |Δ| > σ_seed effect-size criterion (0.434 > 0.406); the paired t-test used throughout (p = 0.169) does not, and the text follows the test.
% NOTE: publication/code/protocol.py declares PROTOCOL.max_epochs = 400, whereas the recorded runs use an 800-epoch cap (set per job by run_matrix.py --max_epochs); the protocol description in §2.3 must say 800.
% EDIT: the backend table (second \label{tab:backend}, duplicate of T4) deleted — see Table~\ref{tab:backend} in tables/all_tables.tex, which carries the per-split timings and the paired p-values.

% src: RESULTS_v1_vs_v2.md §7 (agreement; cheaper implementation; two roughly equal contributions; CPU/float64 caveat; JIT amortised on long uniform runs, interrupted by early stopping);
% src: RESULTS.md §2, ERRATA.md E-19 and PLAYBOOK.md §2.7 (v1 attributed 0.4–0.8 MPa to the backend; the same gap arises between two implementations of one method); MANUSCRIPT.md §4.1.
The implementations are statistically indistinguishable in accuracy (Table~\ref{tab:backend}): $+0.14$~MPa on the random split ($p=0.27$) and $+0.43$~MPa on the extrapolation ($p=0.17$), with the same sign on both. The second gap is of the order of the seed spread: a difference of this size arises between two implementations of one method that differ in their random streams and, here, in the effective physics weight. % EDIT: "differ only in their random streams" contradicted the calibration confound stated above (critic)
Differences of a few tenths of a megapascal between backends are of the same order as the
spread between two implementations of one method and cannot be attributed to the differentiation
engine, so no conclusion about a preferred backend follows from accuracy.
% REV K24: stated as a result of the present comparison rather than as a withdrawal. Where the results are equivalent the cheaper implementation should be used. Here PyTorch completes a full training run 1.65~times faster ($313\pm102$ against $516\pm94$~s pooled over both splits, paired $p=0.002$; per split, $p=0.028$ on \texttt{random} and $p=0.047$ on \texttt{extrap:alpha\_max}, Table~\ref{tab:backend}), for two roughly equal reasons: JAX needs 30\,\% more time per epoch (paired $p=0.005$, $n=10$) and 29\,\% more epochs before early stopping ($p=0.021$; the 800-epoch cap was reached in two JAX runs and one PyTorch run). % src: tables/all_tables.tex tab:backend (wall time: random p = 0.028, extrap:alpha_max p = 0.047, pooled p = 0.002; s/epoch pooled p = 0.005; epochs pooled p = 0.021) % EDIT: p = 0.0009 (unpaired) → paired 0.005; per-split T4 rows cited instead of the pool alone (critic)
The ranking holds for CPU execution in double precision and must not be transferred elsewhere: on a GPU in single precision the just-in-time compilation of JAX is amortised over long uniform runs, whereas here early stopping cuts the runs short.

% =====================================================================
\subsubsection{PINN, MLP and FEM}
\label{sec:res_fem}

% src: runs_v2.csv stage=main, split=random, 5 seeds (per-component RMSE / R², MLP/PINN/VPINN: σ_rr 5.05/5.03/5.07, R² 0.987; σ_θθ 7.04/7.03/6.97, R² 0.990; σ_zz 19.72/19.48/19.53, R² 0.982–0.983; τ_rz 0.363/0.374/0.367, R² 0.778/0.765/0.773 → 0.76–0.78);
% src: QA_DATA_AND_Z.md §4 (axial collapse removes 79 % of τ_rz variance in the 1D track; 2–4 % for normal components); RESULTS.md §5 (8 × 20 grid).
On individual components the families are as close as on the macro figure: on the random split of the one-dimensional track the component RMSE (five seeds) is 5.0--5.1~MPa for $\sigma_{rr}$, 7.0~MPa for $\sigma_{\theta\theta}$, 19.5--19.7~MPa for $\sigma_{zz}$ and 0.36--0.37~MPa for $\tau_{rz}$ for all three families; the low $R^2$ of the shear component (0.76--0.78 against 0.98--0.99 for the normal components; per-component $R^2$ with the held-out standard deviation as denominator) reflects % EDIT: denominator named — csv columns r2_* are local (critic)
the axial collapse of the profile, under which the mean absolute value of $\tau_{rz}$ falls by 79\,\% because the component changes sign along $z$ (Section~\ref{sec:prep-data}). The comparison with the finite-element source is made on the extended track, where $z$ is an input and the full $8\times20$ field is predicted (Section~\ref{sec:models}), because only there can both equilibrium equations be evaluated on the prediction.

% src: git commit 6e7a228 (per-component RMSE/R² and macro on the common 261-case test: MLP 3.52/0.994, 5.87/0.993, 19.45/0.984, 1.25/0.934, 10.33; PINN reduced 3.69/0.993, 6.04/0.993, 20.12/0.983, 1.27/0.932, 10.68; PINN full 3.53/0.994, 5.97/0.993, 20.15/0.982, 1.25/0.935, 10.67);
% src: git commit 89fc7a0 and manuscript/captions.md (dimensionless medians R_r: FEM 0.041, MLP 0.115, PINN reduced 0.064, PINN full 0.062; R_z: 0.091, 0.074, 0.064, 0.047);
% src: git commit 6e7a228, ERRATA.md E-24 and QA_DATA_AND_Z.md §5 (|σ_rr(r=1)|: FEM 3.54, MLP 2.90, PINN reduced 0.78, PINN full 0.56 MPa; ratios ×2.78 and ×1.5 for R_r; BC 4.5 and 6.3 times stricter);
% src: ERRATA.md E-24 "Переобучение подтвердило причину" / commit a79f197: the 1571 → 1498 training pool with the common 261-case test (window-mask cleaning, not the 1756-case union) is documented for the FULL-FORM PINN only.
% TODO: commit 6e7a228 states only that mlp2d and pinn2d_reduced were trained "with the same protocol on the cleaned dataset, common test 261 cases" — the 1498-case pool and the single-seed setting for these two are inferred, not recorded; confirm from the training logs/npz metadata, and consider re-running all three with 5 seeds on the 1756/263 split.
% EDIT: the dataset of this comparison is stated explicitly (critic: 1498 + 261 = 1759 ≠ 1756); numbers checked against ERRATA.md E-24.
This comparison uses an earlier single-seed training on 1498 cases with a common hold-out of 261 cases: the 1848-case set cleaned by the window mask alone (89 cases removed, 1759 in total), on whose random split of 1571/277 the mask removes 73 training and 16 test cases. % src: ERRATA.md E-24 (89 cases with spread > 0.1 mm; 1759 clean; 1571 → 1498 training pool, 261 test; 73 / 16 on the 1848-case random split); code/filters.py unfinished_draw_in_window docstring
The 1756-case set of Section~\ref{sec:prep-data} removes in addition the three runs whose undrawn tail lies outside the window; the other stages of the matrix, including the seed-replicated ablation of Table~\ref{tab:ablation2d}, were run on it. % src: ERRATA.md E-24 (3 runs with the tail outside the window; union 92 → 1756); runs_v2.csv stage twod (1493/263, 1377/379)
The dimensionless residuals of Table~\ref{tab:admissibility} and Fig.~\ref{fig:residual-maps} are the equilibrium residuals $R_r$, $R_z$ of Section~\ref{sec:models} divided by the scales of the loss term \eqref{eq:lphys}, $s_{R_1}=(s_{rr}+s_{\theta\theta})/R_{\mathrm{char}}$ and $s_{R_z}=s_{zz}/L_{\mathrm{char}}+s_{rz}/R_{\mathrm{char}}$, computed once from the FEM data and applied identically to all families, since otherwise each family would be normalised by itself and the rows would not be comparable; a value of 0.1 means that the equation fails to close by 10\,\% of the characteristic magnitude of its own terms. % src: experiments/residual_maps.py residuals() docstring and fig_residuals (scale_r, scale_z; scales from the FEM data, one set for all families; 0.1 = 10 %); captions.md ("Обезразмеривание невязок") % EDIT: definition moved from the caption of tab:admissibility into the text with a reference to eq:lphys (critic; supervisor's instruction)
Medians are taken over the interior grid nodes, where the central difference is second-order accurate; at the edges of the grid it degenerates to a one-sided difference. % src: experiments/residual_maps.py (comment at `inner`); captions.md
The FEM row is the baseline against which every admissibility figure is read. % src: residual_maps.py ("базовая линия — сам МКЭ"); captions.md ("Базовая линия МКЭ")
\begin{table*}[t]
\caption{Accuracy and physical admissibility of the three two-dimensional families on a common hold-out of 261 cases (one training run per family; training pool of 1498 cases, see text). Accuracy: component RMSE (MPa)\,/\,$R^2$ (held-out denominator) and macro-RMSE. Admissibility: median over interior grid nodes and over the test cases of the dimensionless equilibrium residuals $|R_r|$ and $|R_z|$ (definition in the text), and the mean traction-free violation $|\sigma_{rr}(r{=}1)|$. FEM row: the source.} % EDIT: caption made descriptive; definitions moved to the text; R² denominator named — contours.py computes 1 − SS_res/SS_tot over the 261 test cases (critic)
\label{tab:admissibility}
\centering\small
\setlength{\tabcolsep}{4pt}
\begin{tabular}{@{}lccccrrrr@{}}
\toprule
 & \multicolumn{5}{c}{Accuracy, RMSE (MPa)\,/\,$R^2$} & \multicolumn{3}{c}{Admissibility} \\
\cmidrule(lr){2-6}\cmidrule(lr){7-9}
 & $\sigma_{rr}$ & $\sigma_{\theta\theta}$ & $\sigma_{zz}$ & $\tau_{rz}$ & macro & $|R_r|$ & $|R_z|$ & $|\sigma_{rr}(r{=}1)|$, MPa \\
\midrule
FEM (source)       & --           & --           & --            & --           & --    & 0.041 & 0.091 & 3.54 \\
MLP                & 3.52\,/\,0.994 & 5.87\,/\,0.993 & 19.45\,/\,0.984 & 1.25\,/\,0.934 & 10.33 & 0.115 & 0.074 & 2.90 \\
PINN, reduced form & 3.69\,/\,0.993 & 6.04\,/\,0.993 & 20.12\,/\,0.983 & 1.27\,/\,0.932 & 10.68 & 0.064 & 0.064 & 0.78 \\
PINN, full form    & 3.53\,/\,0.994 & 5.97\,/\,0.993 & 20.15\,/\,0.982 & 1.25\,/\,0.935 & 10.67 & 0.062 & 0.047 & 0.56 \\
\bottomrule
\end{tabular}
\end{table*}

% src: Table above; ratios: 0.115/0.041 = 2.8, 0.064/0.041 and 0.062/0.041 = 1.5–1.6, 0.074/0.091 = 0.82, 0.047/0.091 = 0.51, 3.54/0.78 = 4.5, 3.54/0.56 = 6.3 (RESULTS_v1_vs_v2.md §3, §8); macro spread 10.68 − 10.33 = 0.35; residual_maps.py docstring / captions.md (0.1 = equation fails to close by 10 % of its terms).
% Figure labels: fig:compare-fields = figures/paper/fig5_compare_fields.png; fig:residual-maps = figures/paper/fig6_residual_maps.png — the replacement for Figs. 11–12 of the previous version (MANUSCRIPT.md §6.2); captions in manuscript/captions.md. % EDIT: labels renamed fig:compare_fields/fig:residual_maps → fig:compare-fields/fig:residual-maps (one label set for the paper/ figures)
Table~\ref{tab:admissibility} shows where the families part. In accuracy they are equivalent to within 0.35~MPa of macro-RMSE, and their fields are visually alike (Fig.~\ref{fig:compare-fields}, the case with the median error of the full PINN rather than the best). Within each component the colour scale is shared by the FEM reference and all families, so that the panels of one column are directly comparable; in Figs.~\ref{fig:compare-fields} and \ref{fig:residual-maps} $r$ is normalised by the surface radius of the same cross-section and $z$ by the window length, and the radial axis is stretched ($L/R \approx 3.5$). % src: captions.md ("Колорбары и сравнимость", "Нормировка координат"); QA_DATA_AND_Z.md §4 (L/R ≈ 3.5) % EDIT: caption text moved into the body (critic; supervisor's instruction)
In admissibility they are not (Fig.~\ref{fig:residual-maps}). The boundary quantity is evaluated on the material surface, at $r/R = 1$, on the four components of the stress tensor.
% REV K28: the reference to earlier figures removed; the definition that matters is stated directly. % src: ERRATA.md E-01 (R_EXTRAP_FREE_SURFACE = 1.5; condition imposed at 1.5R, measured at r = 1), E-18 (component labels); MANUSCRIPT.md §6.2 % EDIT: sentence added (critic)
The MLP violates radial equilibrium 2.8~times more strongly than the finite-element source it was trained on, both PINNs about 1.5~times; a residual of 0.115 means that the radial equation fails to close by 11.5\,\% of the characteristic magnitude of its own terms. The axial equation is improved only by the full form, the one configuration whose loss contains $R_z$: 0.047 against 0.091 for the source and 0.074 for the MLP. The separation is sharpest at the free surface: the MLP satisfies the traction-free condition no better than the source (2.90 against 3.54~MPa), whereas the reduced and full PINNs satisfy it 4.5 and 6.3~times more strictly than the data they were trained on. Since the finite-element field is affected by nodal extrapolation and averaging, a surrogate that violates equilibrium and the boundary condition by less than its source is, in this specific sense, more physically admissible than the data.

% src: RESULTS.md §1.4 (1D physics audit, 5 seeds; random: |R_1| FEM 402, MLP 1351 ± 146, PINN 749 ± 42, VPINN 1071 ± 26 MPa/m; |σ_rr(r=1)| FEM 3.48, MLP 3.41 ± 0.61, PINN 0.87 ± 0.21, VPINN 0.64 ± 0.14;
%      extrap:alpha_max: |σ_rr(r=1)| FEM 3.35, MLP 3.95 ± 0.82, PINN 1.74 ± 0.32, VPINN 1.80 ± 0.66; extrap:v_max: 31.40 / 33.70 / 28.23 MPa vs FEM 3.10) — all verified.
The ordering is the same with seed statistics on the one-dimensional track. On the random split the median reduced radial residual is 402~MPa/m for the source, $1351\pm146$ for the MLP, $749\pm42$ for the PINN and $1071\pm26$ for the VPINN, and $|\sigma_{rr}(r{=}1)|$ is 3.48, $3.41\pm0.61$, $0.87\pm0.21$ and $0.64\pm0.14$~MPa respectively. The advantage persists with $\alpha=20^\circ$ withheld ($3.95\pm0.82$~MPa for the MLP against $1.74\pm0.32$ and $1.80\pm0.66$, source 3.35) and disappears only where the surrogate itself has failed, on the speed extrapolation (28--34~MPa for every family).

% src: RESULTS.md §5 (2D ablation, 3 seeds: macro-RMSE and full-form residual medians in MPa/m; csv columns eq_res_median = |R_r| full, eq_z_full = |R_z| full) — all 18 entries verified; RESULTS_v1_vs_v2.md §4 (axial −41.4 %/−33.2 % full, −9.2 %/−13.9 % reduced);
% radial percentages computed from RESULTS.md §5 full-form column (626/1040 = −40 %, 553/1040 = −47 %, 739/1426 = −48 %, 686/1426 = −52 %).
% NOTE (critic): RESULTS_v1_vs_v2.md §4/§8 quote −38…−49 % for the radial residual — those are computed from the REDUCED-form residual columns; the text and tab:ablation2d use the full-form columns of RESULTS.md §5 consistently. These residuals are evaluated by the trainer on the 1756-case splits (n_train 1493 / 1377, n_test 263 / 379) and are NOT numerically comparable to the dimensionless medians of tab:admissibility.
% src: QA_DATA_AND_Z.md §4 (dropped axial term = 68 % of retained terms; radial 1.3 %).
\begin{table*}[t]
\caption{Ablation on the extended track ($z$ as input, $8\times20$ grid), three seeds: macro-RMSE (MPa) and median full-form equilibrium residuals (MPa/m) on the held-out cases; in brackets, change relative to the MLP with the same input. The three configurations differ only in the physics terms of the loss.} % EDIT: the sentence on non-comparability with tab:admissibility moved from the caption into the text (critic)
\label{tab:ablation2d}
\centering\small
\begin{tabular}{@{}llrrr@{}}
\toprule
Split & Configuration & Macro-RMSE & $|R_r|$ & $|R_z|$ \\
\midrule
\texttt{random} & MLP, $z$ as input      & $7.81\pm0.22$  & 1040 & 238 \\
                & PINN, reduced form     & $8.05\pm0.44$  & 626 ($-40\,\%$) & 217 ($-9\,\%$) \\
                & PINN, full form        & $8.02\pm0.56$  & 553 ($-47\,\%$) & 140 ($-41\,\%$) \\
\addlinespace
\texttt{extrap:alpha\_max} & MLP, $z$ as input  & $12.70\pm0.03$ & 1426 & 281 \\
                           & PINN, reduced form & $12.18\pm0.34$ & 739 ($-48\,\%$) & 242 ($-14\,\%$) \\
                           & PINN, full form    & $12.21\pm0.51$ & 686 ($-52\,\%$) & 188 ($-33\,\%$) \\
\bottomrule
\end{tabular}
\end{table*}

% src: arithmetic on RESULTS.md §5: 8.02 − 7.81 = +0.21; 12.21 − 12.70 = −0.49.
The ablation in Table~\ref{tab:ablation2d} isolates the axial equation; its residuals are evaluated by the training code in MPa/m on the 1756-case splits and are not numerically comparable to the dimensionless medians of Table~\ref{tab:admissibility}. % src: runs_v2.csv stage twod (eq_res_median, eq_z_full); RESULTS.md §5 % EDIT: moved from the caption (critic)
Both physics-informed forms reduce the radial residual by 40--52\,\%, since both contain the radial equation. The axial residual falls by 33--41\,\% with the full form and by only 9--14\,\% with the reduced one, which drops $\partial\sigma_{zz}/\partial z$, a term amounting to 68\,\% of the retained ones, and therefore cannot enforce axial balance. The gain in axial equilibrium thus comes from the axial term in the loss, not from ``physics in general'', and costs no accuracy: the macro-RMSE changes by $+0.2$~MPa on the random split and $-0.5$~MPa on the extrapolation, within the seed spread.

% src: RESULTS.md §4 (corruption: +250 MPa shift of σ_θθ in a fraction of the training sets, morphology of artefact #439; hold-out clean; macro-RMSE 0/1/5/10 % — all 12 entries verified; degradation ×1.81/×1.64/×1.79, csv 1.813/1.638/1.786);
% src: RESULTS_v1_vs_v2.md §6 (gaps −0.10/−0.48/−1.11/−1.59, p = 0.718/0.100/0.010/0.037; slope −14.20 MPa per unit fraction, p = 0.000363, R² = 0.735) — recomputed from runs_v2.csv stage=corrupt (12 seed-level points): identical.
% NOTE: RESULTS.md §4 quotes p = 0.052 at 10 %: that is an unpaired Welch test (recomputed 0.0518); the seed-paired test used throughout gives 0.037.
% EDIT: the corruption table (tab:corruption, duplicate of T3) deleted — see Table~\ref{tab:corrupt} in tables/all_tables.tex, which also carries the VPINN gap rows and the regression.

The corruption experiment (Table~\ref{tab:corrupt}) yields the strongest and most interpretable effect of the physics term. On clean labels the PINN and the MLP are indistinguishable ($\Delta=-0.10$~MPa, $p=0.72$): the term adds nothing where data and equations already agree. As labels are corrupted the gap opens monotonically, $-0.10\to-0.48\to-1.11\to-1.59$~MPa, and is established at 5 and 10\,\%. A regression of the gap on the corrupted fraction over the twelve seed-level points gives $-14.2$~MPa per unit fraction ($p=0.00036$, $R^2=0.74$): the advantage grows by about 1.4~MPa for every 10\,\% of corrupted labels. At 10\,\% the MLP has degraded by a factor of 1.81 and the PINN by 1.64; the weak form (1.79) is not protected in the same way. The mechanism is the one seen in Table~\ref{tab:admissibility}: the strong-form residual and the traction-free term pull the solution towards an admissible field, and the pull matters exactly when the labels contradict it. Taken together, Tables~\ref{tab:admissibility} and \ref{tab:corrupt} locate the benefit of the physics term not in accuracy against a clean finite-element field, where it is marginal, but in the admissibility of the solution and its resistance to defects of the data source.

% EDIT: Limitations paragraph added at the end of the section (critic) in place of the undefined \ref{sec:limitations}; wording from MANUSCRIPT.md §3, with v2 figures in place of that draft's NRMSE values (0.58 / 2.2–2.4 are v1 and are not used).
\paragraph{Limitations.}
The surrogates are trained and evaluated against a single finite-element model that has not been validated against physical measurements (Section~\ref{sec:prep-data}). This study therefore does not assess how accurately any of the models represents the wire-drawing process; it assesses properties of the surrogates relative to their data source, transferability beyond the training range and robustness to defects of that source, which are measurable without experimental validation and are what is claimed. % src: MANUSCRIPT.md §3
The factorial grid is incomplete and unbalanced: the level $v = 5$~m/min is represented by 122 cases against 554--607 for the other levels, speed coincides with the simulation batch, and 13.2\,\% of the raw cases are bit-identical duplicates of another case differing only in the friction label, so friction cannot be used as an extrapolation axis. % src: MANUSCRIPT.md §3; ERRATA.md E-12 (122 vs 554–607), E-06 (13.2 %); VALIDATION.md §6 (speed = batch)
The shear component $\tau_{rz}$ is essentially not predicted from the one-dimensional representation: its NRMSE (held-out standard deviation as denominator) is 0.47--0.49 on the random hold-out and 1.6--1.9 with the highest speed withheld, against a physical scale of 0.67~MPa mean absolute value, two orders of magnitude below the normal components; macro-averaged figures are therefore dominated by the normal components, and per-component figures are given alongside them. % src: RESULTS.md §1.3б (τ_rz NRMSE random 0.471 / 0.485 / 0.476; extrap:v_max 1.622 / 1.945 / 1.879); PLAYBOOK.md §2.4 (0.67 MPa)
